\section*{Capítulo \ref{ch:g6:inside-the-earth} --- O interior da Terra}

\begin{solution}{exo:g6:inside-the-earth:1}
A \omterm{def:g6:inside-the-earth:layers}{crosta}, o \omterm{def:g6:inside-the-earth:layers}{manto}, o \omterm{def:g6:inside-the-earth:layers}{núcleo externo}, o \omterm{def:g6:inside-the-earth:layers}{núcleo interno}. O \omterm{def:g6:inside-the-earth:layers}{núcleo externo} é
líquido. O \omterm{def:g6:inside-the-earth:layers}{núcleo externo} e o interno são feitos de metal (sobretudo
ferro, com um pouco de níquel); a \omterm{def:g6:inside-the-earth:layers}{crosta} e o \omterm{def:g6:inside-the-earth:layers}{manto} são feitos de \omterm{def:g3:stones-pebbles-sand:rock}{rocha}.
\end{solution}

\begin{solution}{exo:g6:inside-the-earth:2}
(a) Falso: a \omterm{def:g6:inside-the-earth:layers}{crosta} é mais grossa embaixo dos \omterm{def:g3:continents-and-oceans:continent}{continentes} (cerca de 30~km)
que embaixo dos \omterm{def:g2:water-on-earth:ocean}{oceanos} (cerca de 5~km). (b) Falso: o \omterm{def:g6:inside-the-earth:layers}{núcleo} é feito de
metal, uma liga de ferro e de níquel. (c) Verdadeiro. (d) Verdadeiro:
12\,262~m, cerca de 12.3~km.
\end{solution}

\begin{solution}{exo:g6:inside-the-earth:3}
(a) A 2891~km, cerca de 2900~km. (b) A 5150~km. (c) $5150 - 2891 = 2259$:
cerca de 2260~km.
\end{solution}

\begin{solution}{exo:g6:inside-the-earth:4}
``A casca é como a \underline{\omterm{def:g6:inside-the-earth:layers}{crosta}}, a clara é como o \underline{\omterm{def:g6:inside-the-earth:layers}{manto}} e
a gema é como o \underline{\omterm{def:g6:inside-the-earth:layers}{núcleo}}. Mas, ao contrário da gema, o
\underline{\omterm{def:g6:inside-the-earth:layers}{núcleo}} tem duas partes: um \underline{\omterm{def:g6:inside-the-earth:layers}{núcleo externo}} líquido e
um \underline{\omterm{def:g6:inside-the-earth:layers}{núcleo interno}} sólido.''
\end{solution}

\begin{solution}{exo:g6:inside-the-earth:5}
$8.849 \div 6371 \approx 0.00139$, e $0.00139 \times 100 \approx 0.14$: o
Everest tem só cerca de \textbf{0.14~\%} do raio. Vista de longe, a Terra é
uma bola muito lisa: a sua \omterm{def:g1:landscapes-around-us:landform}{montanha} mais alta é uma saliência minúscula.
\end{solution}

\begin{solution}{exo:g6:inside-the-earth:6}
(a) $30 \div 5 = 6$: seis vezes mais grossa. (b) $100 \div 5 = 20$: vinte
vezes mais grossa.
\end{solution}

\begin{solution}{exo:g6:inside-the-earth:7}
(a) $12 + 1 \times 25 = 37~^\circ$C. (b) $12 + 1.5 \times 25 = 12 + 37.5 =
49.5~^\circ$C. (c) $12 + 2 \times 25 = 62~^\circ$C.
\end{solution}

\begin{solution}{exo:g6:inside-the-earth:8}
A subida é de $74 - 14 = 60$ graus em 2~km, então $60 \div 2 =
\textbf{30}$~graus por quilômetro. Sim: 30 fica entre a média do leste (25)
e a média do oeste (34).
\end{solution}

\begin{solution}{exo:g6:inside-the-earth:9}
(a) Logo acima do \omterm{def:g6:inside-the-earth:layers}{núcleo}, cerca de 13.7~km por segundo; logo abaixo do seu
topo, cerca de 8.1~km por segundo: as ondas P ficam mais lentas de repente.
(b) Entre 2891~km e 5150~km, no \omterm{def:g6:inside-the-earth:layers}{núcleo externo}. (c) As ondas S não
conseguem atravessar um líquido: o \omterm{def:g6:inside-the-earth:layers}{núcleo externo} é líquido.
\end{solution}

\begin{solution}{exo:g6:inside-the-earth:10}
(a) $12.3 \div 30 = 0.41$: cerca de \textbf{41~\%} de uma \omterm{def:g6:inside-the-earth:layers}{crosta}
continental média. (b) Não: embaixo de um \omterm{def:g3:continents-and-oceans:continent}{continente}, o \omterm{def:g6:inside-the-earth:layers}{manto} começa
abaixo da \omterm{def:g6:inside-the-earth:layers}{crosta}, a cerca de 30~km de profundidade em média; o furo parou
muito acima dele, na \omterm{def:g6:inside-the-earth:layers}{crosta}.
\end{solution}

\begin{solution}{exo:g6:inside-the-earth:11}
(a) $29 \times 6371 = 184\,759$, e $15 + 184\,759 = 184\,774$: cerca de
185\,000~graus. (b) O laboratório dá cerca de 6000~graus na borda do \omterm{def:g6:inside-the-earth:layers}{núcleo
interno}: o número de Léo é cerca de 30~vezes grande demais
($184\,774 \div 6000 \approx 31$). (c) O ritmo de 29~graus por quilômetro
só foi medido nos primeiros 2~km, mais ou menos. Léo o usa milhares de
quilômetros além das medições; lá no fundo, a temperatura sobe muito mais
devagar. Um ritmo só pode ser usado onde foi medido.
\end{solution}

\begin{solution}{exo:g6:inside-the-earth:12}
(a) Se a \omterm{def:g6:inside-the-earth:layers}{crosta} é mais leve que a média, o material mais fundo tem de ser
mais denso que 5.5, para puxar a média para cima. (b) No \omterm{def:g6:inside-the-earth:layers}{núcleo externo} e
no \omterm{def:g6:inside-the-earth:layers}{núcleo interno} (de cerca de 9.9 a cerca de 13~gramas por centímetro
cúbico); o \omterm{def:g6:inside-the-earth:layers}{manto} mal chega a 5.6 na sua base. (c) Os materiais mais densos
afundam embaixo dos mais leves, como uma pedra afunda na água: o metal
pesado acabou ficando no meio, e as \omterm{def:g3:stones-pebbles-sand:rock}{rochas} mais leves em cima.
\end{solution}

\begin{solution}{pb:g6:inside-the-earth:1}
\textbf{1.} $2 \times 6371 = 12\,742$~km. O modelo tem 2~m de diâmetro, ou
seja, 2000~mm.

\textbf{2.} $12\,742 \div 2000 \approx 6.37$: 1~mm do modelo representa
cerca de 6.37~km da Terra real.

\textbf{3.} $8.849 \div 6.371 \approx 1.39$: o Everest tem cerca de
\textbf{1.4~mm} de altura no modelo, um grão de \omterm{def:g3:stones-pebbles-sand:sand}{areia}.

\textbf{4.} $12.3 \div 6.371 \approx 1.9$~mm. O centro do modelo fica
1000~mm (1~m) abaixo da superfície: o furo é cerca de 520~vezes mais curto
($1000 \div 1.9 \approx 520$).

\textbf{5.} $2891 \div 6.371 \approx 453.8$: o topo do \omterm{def:g6:inside-the-earth:layers}{núcleo externo} é
pintado cerca de \textbf{454~mm} abaixo da superfície.

\textbf{6.} $5150 \div 6.371 \approx 808.3$: cerca de 808~mm abaixo da
superfície. O raio do \omterm{def:g6:inside-the-earth:layers}{núcleo interno} no modelo é de $1000 - 808 = 192$~mm.

\textbf{7.} O \omterm{def:g6:inside-the-earth:layers}{manto}: de cerca de 5~mm a 454~mm, cerca de 449~mm. O \omterm{def:g6:inside-the-earth:layers}{núcleo
externo} ocupa $808 - 454 = 354$~mm, e o \omterm{def:g6:inside-the-earth:layers}{núcleo interno}, 192~mm.

\textbf{8.} $30 \div 6.371 \approx 4.7$~mm: uma faixa fina, que pode ser
pintada com um pincel fino.

\textbf{9.} $15 + 2 \times 29 = 73~^\circ$C. No modelo, 2~km ficam
$2 \div 6.371 \approx 0.31$~mm abaixo da superfície: mal a espessura de
três folhas de papel.

\textbf{10.} $4 \div 6.371 \approx 0.63$~mm, contra cerca de 1.9~mm para o
furo: a mina é mais rasa, por cerca de $1.9 - 0.6 = 1.3$~mm.

\textbf{11.} Derreter ou não depende da temperatura e da pressão. O \omterm{def:g6:inside-the-earth:layers}{núcleo
interno} é o mais apertado de todos, por cerca de 3.3~milhões de vezes a
pressão do ar: o ferro apertado só derrete a uma temperatura mais alta,
então o ferro do \omterm{def:g6:inside-the-earth:layers}{núcleo interno} continua sólido, embora seja mais quente
que o \omterm{def:g6:inside-the-earth:layers}{núcleo externo}.

\textbf{12.} $5 \div 6.371 \approx 0.785$: a \omterm{def:g6:inside-the-earth:layers}{crosta} oceânica deve ser
pintada com só \textbf{0.78~mm} de espessura no modelo de 2~metros, a
espessura de cerca de oito folhas de papel ($0.78 \div 0.1 \approx 8$).
\end{solution}
